\documentclass[a4paper,11pt]{ctexart}
\usepackage{amsmath, amssymb, amsfonts}
\usepackage{geometry}
\geometry{left=1.5cm, right=1.5cm, top=2.0cm, bottom=2.0cm, headsep=0.3cm, footskip=1cm}
\usepackage{listings}
\usepackage{xcolor}
\usepackage{tikz}
\usepackage{pgfplots}
\usepackage{hyperref}
\usepackage{fancyhdr}
\usepackage{tcolorbox}
\tcbuselibrary{breakable, skins}
\usepackage{array}
\usepackage{booktabs}
\usepackage[shortlabels]{enumitem}
\usepackage{needspace}
\usepackage{multirow}

\AtBeginDocument{\hypersetup{colorlinks=true, linkcolor=blue, filecolor=magenta, urlcolor=cyan}}
\setlist{nosep, itemsep=0.8ex, parsep=0.1em}
\setlist[itemize]{leftmargin=1.8em}
\setlist[enumerate]{leftmargin=1.8em}

\definecolor{codegreen}{rgb}{0,0.6,0}
\definecolor{codegray}{rgb}{0.5,0.5,0.5}
\definecolor{codepurple}{rgb}{0.58,0,0.82}
\definecolor{codebg}{rgb}{0.98,0.98,0.98}
\definecolor{tipsblue}{rgb}{0.85,0.93,1.0}
\definecolor{highlightyellow}{rgb}{1.0,0.97,0.8}

\lstdefinestyle{mystyle}{backgroundcolor=\color{codebg},commentstyle=\color{codegreen},keywordstyle=\color{magenta}\bfseries,numberstyle=\tiny\color{codegray},stringstyle=\color{codepurple},basicstyle=\ttfamily\small,breakatwhitespace=false,breaklines=true,captionpos=b,keepspaces=true,numbers=left,numbersep=6pt,showspaces=false,showstringspaces=false,showtabs=false,tabsize=2,language=Python,frame=single,framerule=0.5pt,rulecolor=\color{gray!40}}
\lstset{style=mystyle}

\newtcolorbox{problembox}[1][]{colback=white,colframe=blue!60!black,title=\textbf{#1},fonttitle=\bfseries\small,boxrule=1pt,arc=3pt,outer arc=3pt,coltitle=black,before skip=10pt,after skip=8pt}
\newtcolorbox{solutionbox}[1][]{enhanced,breakable,colback=green!3!white,colframe=green!50!black,title=\textbf{#1},fonttitle=\bfseries\small,boxrule=1pt,arc=3pt,outer arc=3pt,coltitle=black,before skip=8pt,after skip=10pt}
\newtcolorbox{metaphorbox}[1][]{colback=orange!5!white,colframe=orange!70!black,title=\textbf{#1},fonttitle=\bfseries\small,boxrule=1pt,arc=3pt,outer arc=3pt,coltitle=black,before skip=8pt,after skip=8pt}
\newtcolorbox{beginnerbox}[1][]{colback=tipsblue,colframe=blue!50!black,title=\textbf{#1},fonttitle=\bfseries\small,boxrule=1pt,arc=3pt,outer arc=3pt,coltitle=black,before skip=8pt,after skip=8pt}
\newtcolorbox{appbox}[1][]{colback=green!3!white,colframe=green!60!black,title=\textbf{#1},fonttitle=\bfseries\small,boxrule=1pt,arc=3pt,outer arc=3pt,coltitle=black,before skip=8pt,after skip=10pt}
\newtcolorbox{keybox}[1][]{colback=highlightyellow,colframe=orange!80!black,title=\textbf{#1},fonttitle=\bfseries\small,boxrule=1pt,arc=3pt,outer arc=3pt,coltitle=black,before skip=8pt,after skip=8pt}

\pagestyle{fancy}
\fancyhf{}
\fancyhead[R]{机器学习-第4章}
\fancyhead[L]{模型评估和模型选择 · 练习题}
\fancyfoot[C]{\thepage}
\addtolength{\headheight}{2pt}

\parskip=2pt plus 1pt minus 0.5pt
\linespread{1.35}
\raggedbottom
\widowpenalty=10000
\clubpenalty=10000

\title{\textbf{机器学习\\第4章\quad 模型评估和模型选择}\\[0.5em]
{\large\color{gray}——如何知道模型学得好不好？如何让它学得更好？}}
\date{\today}

\begin{document}
\maketitle

\begin{beginnerbox}[本章题目导览：你将挑战哪些核心问题？]
    本章是机器学习的"质量控制"核心章节——没有评估，就不知道模型在现实中会不会失败。本章 8 道题覆盖从损失函数到梯度下降的完整评估与优化体系：
    \begin{itemize}
        \item \textbf{Q4-01（损失函数）}：四种损失函数的公式、性质与适用场景——为什么不同任务用不同的损失？
        \item \textbf{Q4-02（欠拟合/过拟合）}：识别、诊断、解决——偏差-方差权衡是什么？
        \item \textbf{Q4-03（正则化计算）}：L1/L2/ElasticNet 的公式、几何直觉与超参数作用，手动计算正则化项。
        \item \textbf{Q4-04（交叉验证）}：Hold-Out、k 折、LOO 三种策略的优缺点与适用场景。
        \item \textbf{Q4-05（梯度下降推导）}：BGD/SGD/MBGD 的对比，以及梯度下降一步迭代的手动计算。
        \item \textbf{Q4-06（回归评估指标）}：MAE/MSE/RMSE/$R^2$ 的公式计算与含义辨析。
        \item \textbf{Q4-07（分类评估指标）}：混淆矩阵、准确率/精确率/召回率/F1 的计算，类别不平衡下的指标选择。
        \item \textbf{Q4-08（ROC/AUC）}：ROC 曲线的绘制原理、AUC 的含义，以及 ROC vs PR 曲线的选择。
    \end{itemize}
\end{beginnerbox}

\section{第 4 章\quad 模型评估和模型选择}

\begin{metaphorbox}[先建立一个总感觉：评估是什么？为什么重要？]
    模型训练完之后，我们面临两个关键问题：\textbf{它学得好吗？}（评估）\textbf{怎么让它学得更好？}（选择与优化）。\\
    没有评估就没有方向——就像考试不公布分数，学生不知道哪里错了，怎么改进？\\
    评估的难点在于：我们真正关心的是模型在\textbf{未知新数据}上的表现，但手里只有有限的历史数据。所有评估方法（交叉验证、测试集）都是在用有限数据\textbf{估计}这个"真实表现"。
\end{metaphorbox}

%% ============================================================
\Needspace{5\baselineskip}
\begin{problembox}[Q4-01\quad 损失函数：四种损失的公式与适用场景]
    \textbf{题目描述：}\\
    损失函数量化了模型预测与真实值之间的差距，是策略的核心。不同任务对损失函数有不同的要求。

    \vspace{0.3cm}
    \textbf{问题：}
    \begin{enumerate}[(1)]
        \item 填写下表，对比四种常见损失函数：

        \begin{center}
        \begin{tabular}{p{2.8cm} p{3.5cm} p{3cm} p{3cm}}
            \toprule
            损失函数 & 公式 & 适用任务 & 主要特点 \\
            \midrule
            0-1损失 & $L = \mathbb{1}[\hat{y} \neq y]$ & \underline{\hspace{2.5cm}} & \underline{\hspace{2.5cm}} \\[1.2em]
            平方损失（MSE） & \underline{\hspace{3cm}} & \underline{\hspace{2.5cm}} & 对异常值敏感 \\[1.2em]
            绝对损失（MAE） & \underline{\hspace{3cm}} & \underline{\hspace{2.5cm}} & \underline{\hspace{2.5cm}} \\[1.2em]
            对数似然损失 & $-[y\log\hat{p}+(1-y)\log(1-\hat{p})]$ & \underline{\hspace{2.5cm}} & \underline{\hspace{2.5cm}} \\
            \bottomrule
        \end{tabular}
        \end{center}

        \item 0-1 损失函数为什么\textbf{不能用梯度下降}来优化？实际分类任务中为什么用对数损失来替代它？

        \item 平方损失和绝对损失都可以用于回归，但有一个关键差异：异常值（Outlier）存在时，哪种损失函数更鲁棒？从数学角度解释原因。

        \item 写出多分类任务的交叉熵损失公式（Softmax Cross-Entropy），并说明其与二分类对数损失的关系。
    \end{enumerate}
    {\color{gray}\small\textit{来源溯源：[文档第 4 章 4.1 损失函数]}}
\end{problembox}

\begin{solutionbox}[参考答案与详细解析]
    \textbf{(1) 四种损失对比：}
    \begin{center}
    \begin{tabular}{p{2.8cm} p{3.5cm} p{3cm} p{3cm}}
        \toprule
        损失函数 & 公式 & 适用任务 & 主要特点 \\
        \midrule
        0-1损失 & $\mathbb{1}[\hat{y}\neq y]$ & 分类 & 不可微，无法梯度优化 \\[0.8em]
        平方损失 & $(\hat{y}-y)^2$ & 回归 & 处处可微，对异常值极敏感 \\[0.8em]
        绝对损失 & $|\hat{y}-y|$ & 回归 & 在0处不可微，对异常值鲁棒 \\[0.8em]
        对数似然损失 & $-[y\log\hat{p}+(1-y)\log(1-\hat{p})]$ & 二分类 & 概率化损失，凸函数，梯度平滑 \\
        \bottomrule
    \end{tabular}
    \end{center}

    \textbf{(2) 0-1 损失不可梯度优化：}

    0-1 损失是阶跃函数：在正确预测时梯度为 0，在错误预测时梯度无定义（不连续跳变）——无法计算有意义的梯度方向，梯度下降无从下手。对数损失是 0-1 损失的\textbf{凸连续可微上界}：当预测正确时损失接近 0，预测错误时损失趋向无穷，惩罚了"高置信度的错误预测"，同时处处可微，梯度可以平滑地传播。

    \textbf{(3) 对异常值的鲁棒性：}

    \textbf{MAE（绝对损失）更鲁棒}。平方损失对误差做平方放大——一个误差为 10 的异常点产生损失 100，贡献的梯度与误差成正比，会严重"扭曲"模型参数去迁就这个异常点。MAE 的梯度是恒定的 $\pm 1$（与误差大小无关），异常点的影响被限制在固定范围内，不会因为误差大而主导整个优化方向。

    \textbf{(4) 多分类交叉熵：}
    \[
    L = -\sum_{c=1}^C y_c \log \hat{p}_c
    \]
    其中 $\hat{p}_c = \text{softmax}(z_c)$，$y_c$ 是 one-hot 标签。当 $C=2$ 时，$y_1 = y$，$y_2 = 1-y$，退化为二分类对数损失，两者是统一框架下的特例关系。
\end{solutionbox}

%% ============================================================
\Needspace{5\baselineskip}
\begin{problembox}[Q4-02\quad 欠拟合与过拟合：诊断、原因与解决]
    \textbf{题目描述：}\\
    欠拟合和过拟合是机器学习中最核心的两个问题，理解它们的本质是选择合适模型和正则化策略的基础。

    \vspace{0.3cm}
    \textbf{问题：}
    \begin{enumerate}[(1)]
        \item 解释\textbf{偏差（Bias）}和\textbf{方差（Variance）}的含义，并填写下表：

        \begin{center}
        \begin{tabular}{p{2.5cm} p{3cm} p{3cm} p{3.5cm}}
            \toprule
            & 欠拟合 & 过拟合 & 理想模型 \\
            \midrule
            偏差 & \underline{\hspace{2.5cm}} & \underline{\hspace{2.5cm}} & 低 \\
            方差 & \underline{\hspace{2.5cm}} & \underline{\hspace{2.5cm}} & 低 \\
            训练集误差 & \underline{\hspace{2.5cm}} & \underline{\hspace{2.5cm}} & 低 \\
            测试集误差 & \underline{\hspace{2.5cm}} & \underline{\hspace{2.5cm}} & 低（接近训练误差）\\
            产生原因 & \underline{\hspace{2.5cm}} & \underline{\hspace{2.5cm}} & — \\
            \bottomrule
        \end{tabular}
        \end{center}

        \item \textbf{诊断方法}：如何通过"训练误差"和"验证误差"的大小关系来判断模型处于欠拟合还是过拟合状态？

        \item \textbf{解决欠拟合的方法}（至少 3 种）；\textbf{解决过拟合的方法}（至少 4 种）。

        \item 学习曲线（Learning Curve）是一种常用的诊断工具，横轴是训练样本数量，纵轴是误差。描述欠拟合和过拟合模型各自的学习曲线形状特征。
    \end{enumerate}
    {\color{gray}\small\textit{来源溯源：[文档第 4 章 4.3 欠拟合与过拟合]}}
\end{problembox}

\begin{solutionbox}[参考答案与详细解析]
    \textbf{(1) 偏差-方差表格：}
    \begin{center}
    \begin{tabular}{p{2.5cm} p{3cm} p{3cm} p{3.5cm}}
        \toprule
        & 欠拟合 & 过拟合 & 理想模型 \\
        \midrule
        偏差 & 高（系统性错误） & 低 & 低 \\
        方差 & 低 & 高（对数据敏感） & 低 \\
        训练集误差 & 高 & 很低（近乎为0） & 低 \\
        测试集误差 & 高（与训练误差相近） & 很高（远大于训练误差） & 低 \\
        产生原因 & 模型太简单/特征不足 & 模型太复杂/数据太少 & — \\
        \bottomrule
    \end{tabular}
    \end{center}

    \textbf{(2) 诊断方法：}
    \begin{itemize}
        \item 训练误差高 + 验证误差高 且两者接近 → \textbf{欠拟合}（高偏差）
        \item 训练误差低 + 验证误差远高于训练误差 → \textbf{过拟合}（高方差）
        \item 训练误差低 + 验证误差也低且接近 → \textbf{理想状态}
    \end{itemize}

    \textbf{(3) 解决方案：}
    \begin{itemize}
        \item \textbf{欠拟合}：增加模型复杂度（更深的树、更多层神经网络）；增加特征（特征工程构造新特征）；减少正则化强度；训练更长时间（增加迭代轮数）。
        \item \textbf{过拟合}：增加训练数据量；减少模型复杂度（降低树深/减少层数）；正则化（L1/L2/Dropout）；数据增强；早停（Early Stopping）；交叉验证选超参数。
    \end{itemize}

    \textbf{(4) 学习曲线特征：}

    \begin{keybox}[学习曲线诊断]
        \textbf{欠拟合}：训练误差和验证误差都\textbf{很高}，随数据量增加两者缓慢下降但最终趋于一个较高的值，且两者间隙小——增加数据无法根本解决问题，需要改进模型。\\
        \textbf{过拟合}：训练误差\textbf{很低}，验证误差\textbf{很高}，两者间隙很大。随数据量增加，验证误差下降，训练误差上升，两者逐渐靠拢——增加数据有效，或添加正则化。
    \end{keybox}
\end{solutionbox}

%% ============================================================
\Needspace{5\baselineskip}
\begin{problembox}[Q4-03\quad 正则化：L1/L2/ElasticNet 的公式与手算]
    \textbf{题目描述：}\\
    正则化通过在损失函数中添加惩罚项来限制模型复杂度，防止过拟合。三种主要正则化方法各有特点。

    \vspace{0.2cm}
    \textbf{问题：}
    \begin{enumerate}[(1)]
        \item 写出加入正则化后的目标函数一般形式（以线性模型为例），并说明正则化系数 $\lambda$ 的作用。

        \item 设模型参数 $\mathbf{w} = [2, -1, 0.5, -3]$，$\lambda = 0.1$，分别计算：
        \begin{enumerate}[a.]
            \item L1 正则化项 $R_1 = \lambda\|\mathbf{w}\|_1$
            \item L2 正则化项 $R_2 = \lambda\|\mathbf{w}\|_2^2$
            \item ElasticNet 正则化项（$\alpha=0.5$，即 L1 和 L2 各半权重）
        \end{enumerate}

        \item 填写下表，对比三种正则化方法：
        \begin{center}
        \begin{tabular}{p{2.5cm} p{3.5cm} p{3.5cm} p{3cm}}
            \toprule
            方法 & 数学形式 & 对参数的作用 & 典型使用场景 \\
            \midrule
            L1（Lasso） & $\lambda\sum|w_j|$ & 将部分参数\underline{\hspace{1.5cm}} & \underline{\hspace{2.5cm}} \\[1.2em]
            L2（Ridge） & \underline{\hspace{3cm}} & 将参数\underline{\hspace{2cm}} & \underline{\hspace{2.5cm}} \\[1.2em]
            ElasticNet & \underline{\hspace{3cm}} & \underline{\hspace{3cm}} & \underline{\hspace{2.5cm}} \\
            \bottomrule
        \end{tabular}
        \end{center}

        \item 参数 $\lambda$ 值越大，模型会发生什么变化（过拟合方向还是欠拟合方向）？$\lambda=0$ 和 $\lambda\to\infty$ 各对应什么极端情况？
    \end{enumerate}
    {\color{gray}\small\textit{来源溯源：[文档第 4 章 4.4 正则化]}}
\end{problembox}

\begin{solutionbox}[参考答案与详细解析]
    \textbf{(1) 正则化目标函数：}
    \[
    J(\mathbf{w}) = \underbrace{\frac{1}{n}\sum_{i=1}^n L(f(x_i), y_i)}_{\text{经验损失}} + \underbrace{\lambda R(\mathbf{w})}_{\text{正则化项}}
    \]
    $\lambda > 0$ 控制正则化强度：$\lambda$ 越大，对参数幅度的惩罚越重，模型参数被压缩得越小，模型越简单（偏向欠拟合）；$\lambda$ 越小，正则化影响越弱，模型可以更自由地拟合数据。

    \textbf{(2) 手动计算（$\mathbf{w}=[2,-1,0.5,-3], \lambda=0.1$）：}
    \begin{enumerate}[a.]
        \item $R_1 = 0.1 \times (|2|+|-1|+|0.5|+|-3|) = 0.1 \times 6.5 = \mathbf{0.65}$
        \item $R_2 = 0.1 \times (4+1+0.25+9) = 0.1 \times 14.25 = \mathbf{1.425}$
        \item ElasticNet（$\alpha=0.5$）：$R = \lambda[\alpha\|\mathbf{w}\|_1 + (1-\alpha)\|\mathbf{w}\|_2^2] = 0.1[0.5\times6.5 + 0.5\times14.25] = 0.1\times10.375 = \mathbf{1.0375}$
    \end{enumerate}

    \textbf{(3) 三种正则化对比：}
    \begin{center}
    \begin{tabular}{p{2.5cm} p{3.5cm} p{3.5cm} p{3cm}}
        \toprule
        方法 & 数学形式 & 对参数的作用 & 典型场景 \\
        \midrule
        L1 & $\lambda\sum|w_j|$ & 将部分参数\textbf{压缩为精确的零}（稀疏） & 特征选择、高维稀疏场景 \\[0.8em]
        L2 & $\lambda\sum w_j^2$ & 将参数\textbf{均匀压小但不为零} & 多重共线性、大多数回归场景 \\[0.8em]
        ElasticNet & $\lambda[\alpha\|\mathbf{w}\|_1+(1-\alpha)\|\mathbf{w}\|_2^2]$ & 兼具稀疏性和稳定性 & 特征间高度相关且数量多的场景 \\
        \bottomrule
    \end{tabular}
    \end{center}

    \textbf{(4) $\lambda$ 的极端情况：}

    $\lambda$ 越大，正则化项主导目标函数，模型参数被强行压向 0，拟合能力减弱，\textbf{偏向欠拟合}。$\lambda=0$：无正则化，模型可自由拟合，\textbf{退化为无约束的普通回归}（过拟合风险最高）。$\lambda\to\infty$：正则化项极大，所有参数趋向 0（L1）或趋向 0（L2），模型\textbf{退化为只输出均值的常数预测}（严重欠拟合）。
\end{solutionbox}

%% ============================================================
\Needspace{5\baselineskip}
\begin{problembox}[Q4-04\quad 交叉验证：Hold-Out、k 折、LOO 对比]
    \textbf{题目描述：}\\
    交叉验证是通过合理划分数据来可靠估计模型泛化误差的方法，有三种主要策略。

    \vspace{0.2cm}
    \textbf{问题：}
    \begin{enumerate}[(1)]
        \item 填写下表对比三种交叉验证策略：
        \begin{center}
        \begin{tabular}{p{2.5cm} p{3cm} p{3.5cm} p{3.5cm}}
            \toprule
            方法 & 核心思路 & 优点 & 缺点/适用条件 \\
            \midrule
            Hold-Out & 一次性划分训练/验证集 & \underline{\hspace{2.5cm}} & \underline{\hspace{3cm}} \\[1.2em]
            k折交叉验证 & \underline{\hspace{2.5cm}} & \underline{\hspace{2.5cm}} & \underline{\hspace{3cm}} \\[1.2em]
            LOO & \underline{\hspace{2.5cm}} & \underline{\hspace{2.5cm}} & \underline{\hspace{3cm}} \\
            \bottomrule
        \end{tabular}
        \end{center}

        \item k 折交叉验证中，$k$ 一般取多少？$k$ 取得太大（如 $k=n$，退化为 LOO）和太小（如 $k=2$）分别有什么问题？

        \item 写出 sklearn 实现 5 折交叉验证的核心代码，并说明 \texttt{cross\_val\_score} 返回的是什么。

        \item 分层 k 折交叉验证（Stratified k-Fold）和普通 k 折有什么区别？在什么情况下必须用分层交叉验证？
    \end{enumerate}
    {\color{gray}\small\textit{来源溯源：[文档第 4 章 4.5 交叉验证]}}
\end{problembox}

\begin{solutionbox}[参考答案与详细解析]
    \textbf{(1) 三种策略对比：}
    \begin{center}
    \begin{tabular}{p{2.5cm} p{3cm} p{3.5cm} p{3.5cm}}
        \toprule
        方法 & 核心思路 & 优点 & 缺点/适用条件 \\
        \midrule
        Hold-Out & 固定比例划分，训练/验证各用一次 & 速度快，实现简单 & 评估结果依赖随机划分，方差大；需要数据量充足 \\[0.8em]
        k折CV & 数据分 k 份，轮流用 1 份验证、k-1 份训练，取均值 & 充分利用数据，评估更稳定可靠 & 训练 k 次，计算代价高；k 的选择影响偏差-方差 \\[0.8em]
        LOO & 每次留 1 个样本验证，其余全部训练 & 最充分利用数据，接近无偏估计 & 计算代价极高（训练 n 次）；评估方差高，只适合极小数据集 \\
        \bottomrule
    \end{tabular}
    \end{center}

    \textbf{(2) $k$ 值的选择：}

    通常取 $k=5$ 或 $k=10$（工业经验最佳平衡点）。$k$ 太大（接近 LOO）：每次训练集更多、验证集更小，偏差低但评估方差高，计算代价极大。$k$ 太小（如 $k=2$）：每次训练集只有一半数据，模型欠拟合，评估偏差大。

    \textbf{(3) sklearn 5 折交叉验证代码：}
    \begin{lstlisting}[language=Python]
from sklearn.model_selection import cross_val_score
from sklearn.linear_model import LogisticRegression

model = LogisticRegression()
scores = cross_val_score(model, X, y, cv=5, scoring='accuracy')
print(f"5折CV各折准确率: {scores}")
print(f"均值 ± 标准差: {scores.mean():.3f} ± {scores.std():.3f}")
    \end{lstlisting}
    \texttt{cross\_val\_score} 返回一个长度为 $k$ 的数组，每个元素是对应折验证集上的评估分数。取均值作为模型的最终估计性能，标准差衡量估计的稳定性。

    \textbf{(4) 分层 k 折：}

    分层 k 折在划分时\textbf{保持每折中各类别比例与原始数据一致}（如原始数据正负比例 1:9，每折也保持 1:9）。当数据类别极不平衡时必须使用——若不分层，某折可能完全没有少数类样本，导致该折评估失真（如准确率虚高）。
\end{solutionbox}

%% ============================================================
\Needspace{5\baselineskip}
\begin{problembox}[Q4-05\quad 梯度下降：推导 + 手动计算一步迭代]
    \textbf{题目描述：}\\
    梯度下降是最核心的参数优化算法，三种变体各有取舍。本题通过手动计算加深理解。

    \vspace{0.2cm}
    \textbf{问题：}
    \begin{enumerate}[(1)]
        \item 梯度下降的参数更新公式是 $\theta \leftarrow \theta - \eta \nabla_\theta J(\theta)$。解释每个符号的含义，并说明为什么是"负梯度"方向。

        \item 设损失函数为 $J(w) = (w - 3)^2$，当前参数 $w_0 = 0$，学习率 $\eta = 0.1$。手动执行 3 步梯度下降，填写下表：

        \begin{center}
        \begin{tabular}{ccccc}
            \toprule
            步骤 & 当前 $w$ & 梯度 $\frac{dJ}{dw}$ & 更新量 $\eta \cdot \text{梯度}$ & 新 $w$ \\
            \midrule
            1 & 0 & \underline{\quad} & \underline{\quad} & \underline{\quad} \\
            2 & \underline{\quad} & \underline{\quad} & \underline{\quad} & \underline{\quad} \\
            3 & \underline{\quad} & \underline{\quad} & \underline{\quad} & \underline{\quad} \\
            \bottomrule
        \end{tabular}
        \end{center}

        \item 填写下表，对比三种梯度下降变体：
        \begin{center}
        \begin{tabular}{p{2.8cm} p{2.5cm} p{3cm} p{3.5cm}}
            \toprule
            变体 & 每次用多少数据 & 收敛稳定性 & 适用场景 \\
            \midrule
            BGD & \underline{\hspace{2cm}} & \underline{\hspace{2.5cm}} & \underline{\hspace{3cm}} \\[1.2em]
            SGD & \underline{\hspace{2cm}} & \underline{\hspace{2.5cm}} & \underline{\hspace{3cm}} \\[1.2em]
            MBGD & \underline{\hspace{2cm}} & \underline{\hspace{2.5cm}} & \underline{\hspace{3cm}} \\
            \bottomrule
        \end{tabular}
        \end{center}

        \item 学习率 $\eta$ 过大和过小分别会导致什么问题？如何选择合适的学习率？
    \end{enumerate}
    {\color{gray}\small\textit{来源溯源：[文档第 4 章 4.6.2 梯度下降法]}}
\end{problembox}

\begin{solutionbox}[参考答案与详细解析]
    \textbf{(1) 公式符号含义：}

    $\theta$：模型参数；$\eta$：学习率（步长大小）；$\nabla_\theta J(\theta)$：损失函数对参数的梯度，指向损失\textbf{上升最快}的方向。取\textbf{负梯度}方向（$-\nabla J$）是因为我们要\textbf{最小化}损失，沿梯度反方向走能最快减小损失值——爬山时反向走就是下山。

    \textbf{(2) 手动三步迭代（$J(w)=(w-3)^2$，梯度 $\frac{dJ}{dw}=2(w-3)$）：}

    \begin{center}
    \begin{tabular}{ccccc}
        \toprule
        步骤 & 当前 $w$ & 梯度 $2(w-3)$ & 更新量 $0.1\times$梯度 & 新 $w$ \\
        \midrule
        1 & 0 & $-6$ & $-0.6$ & $0.6$ \\
        2 & $0.6$ & $-4.8$ & $-0.48$ & $1.08$ \\
        3 & $1.08$ & $-3.84$ & $-0.384$ & $1.464$ \\
        \bottomrule
    \end{tabular}
    \end{center}
    最优解为 $w^*=3$，可见每步都在向 3 靠近（以 0.8 的比例收缩剩余距离）。

    \textbf{(3) 三种变体对比：}
    \begin{center}
    \begin{tabular}{p{2.8cm} p{2.5cm} p{3cm} p{3.5cm}}
        \toprule
        变体 & 每次用多少数据 & 收敛稳定性 & 适用场景 \\
        \midrule
        BGD & 全量数据 & 稳定，单调收敛 & 数据量小，内存足够；凸优化问题 \\[0.8em]
        SGD & 1 个样本 & 不稳定，震荡 & 在线学习；超大规模数据集 \\[0.8em]
        MBGD & mini-batch（32/64/128） & 较稳定，可并行 & \textbf{深度学习最常用}；GPU 并行加速 \\
        \bottomrule
    \end{tabular}
    \end{center}

    \textbf{(4) 学习率选择：}

    $\eta$ 过大：步长太大，可能越过最优点，在谷底两侧来回震荡甚至发散（损失不降反升）。$\eta$ 过小：收敛极慢，需要大量迭代；容易陷入局部极小值或鞍点。\textbf{选择策略}：使用学习率调度（初始较大，随训练衰减）；网格搜索或学习率预热（Warm-up）；使用自适应学习率优化器（Adam、AdaGrad）自动调节每个参数的学习率。
\end{solutionbox}

%% ============================================================
\Needspace{5\baselineskip}
\begin{problembox}[Q4-06\quad 回归评估指标：MAE/MSE/RMSE/$R^2$ 手动计算]
    \textbf{题目描述：}\\
    设某回归模型对 5 个样本的预测结果如下：

    \begin{center}
    \begin{tabular}{cccc}
        \toprule
        样本 & 真实值 $y_i$ & 预测值 $\hat{y}_i$ & 误差 $e_i = \hat{y}_i - y_i$ \\
        \midrule
        1 & 10 & 12 & $+2$ \\
        2 & 20 & 18 & $-2$ \\
        3 & 15 & 15 & $0$ \\
        4 & 25 & 30 & $+5$ \\
        5 & 30 & 27 & $-3$ \\
        \bottomrule
    \end{tabular}
    \end{center}

    \textbf{问题：}
    \begin{enumerate}[(1)]
        \item 计算 \textbf{MAE}（平均绝对误差）。
        \item 计算 \textbf{MSE}（均方误差）和 \textbf{RMSE}（均方根误差）。
        \item 计算 $\mathbf{R^2}$（决定系数）。（提示：$\bar{y} = 20$，$\text{SS\_tot} = \sum(y_i-\bar{y})^2$，$\text{SS\_res} = \sum(y_i-\hat{y}_i)^2$）
        \item 填写下表，说明四个指标各自的含义与优缺点：
        \begin{center}
        \begin{tabular}{p{2cm} p{4cm} p{4.5cm} p{2cm}}
            \toprule
            指标 & 含义 & 优缺点 & 取值范围 \\
            \midrule
            MAE & \underline{\hspace{3.5cm}} & \underline{\hspace{4cm}} & \underline{\quad} \\[1.2em]
            RMSE & \underline{\hspace{3.5cm}} & \underline{\hspace{4cm}} & \underline{\quad} \\[1.2em]
            $R^2$ & \underline{\hspace{3.5cm}} & \underline{\hspace{4cm}} & $(-\infty, 1]$ \\
            \bottomrule
        \end{tabular}
        \end{center}
    \end{enumerate}
    {\color{gray}\small\textit{来源溯源：[文档第 4 章 4.7.1 回归模型评价指标]}}
\end{problembox}

\begin{solutionbox}[参考答案与详细解析]
    \textbf{(1) MAE：}
    \[
    \text{MAE} = \frac{1}{5}(|2|+|-2|+|0|+|5|+|-3|) = \frac{2+2+0+5+3}{5} = \frac{12}{5} = \mathbf{2.4}
    \]

    \textbf{(2) MSE 和 RMSE：}
    \[
    \text{MSE} = \frac{1}{5}(4+4+0+25+9) = \frac{42}{5} = \mathbf{8.4}
    \]
    \[
    \text{RMSE} = \sqrt{8.4} \approx \mathbf{2.898}
    \]

    \textbf{(3) $R^2$：}

    $\bar{y}=20$，各真实值与均值之差：$(10-20)^2+(20-20)^2+(15-20)^2+(25-20)^2+(30-20)^2 = 100+0+25+25+100=250$

    $\text{SS\_tot} = 250$，$\text{SS\_res} = 4+4+0+25+9 = 42$
    \[
    R^2 = 1 - \frac{\text{SS\_res}}{\text{SS\_tot}} = 1 - \frac{42}{250} = 1 - 0.168 = \mathbf{0.832}
    \]
    说明该模型解释了目标变量 83.2\% 的方差。

    \textbf{(4) 指标含义与优缺点：}
    \begin{center}
    \begin{tabular}{p{2cm} p{4cm} p{4.5cm} p{2cm}}
        \toprule
        指标 & 含义 & 优缺点 & 取值范围 \\
        \midrule
        MAE & 预测误差绝对值的均值，与目标变量同量纲 & 直观易解释，对异常值鲁棒；但不可微（绝对值在0处） & $[0,+\infty)$ \\[0.8em]
        RMSE & 预测误差平方均值的开方，与目标变量同量纲 & 对大误差惩罚更重，可微；对异常值敏感 & $[0,+\infty)$ \\[0.8em]
        $R^2$ & 模型解释了目标变量多少比例的方差（1=完美，0=仅预测均值） & 无量纲，便于不同数据集比较；可为负（模型差于均值预测） & $(-\infty,1]$ \\
        \bottomrule
    \end{tabular}
    \end{center}
\end{solutionbox}

%% ============================================================
\Needspace{5\baselineskip}
\begin{problembox}[Q4-07\quad 分类评估指标：混淆矩阵到 F1 分数]
    \textbf{题目描述：}\\
    某二分类模型（预测是否患病）在测试集上的预测结果如下：

    \begin{center}
    \begin{tabular}{ccc}
        \toprule
        & 预测为正（患病） & 预测为负（未患病） \\
        \midrule
        实际为正（患病） & TP = 40 & FN = 10 \\
        实际为负（未患病） & FP = 20 & TN = 130 \\
        \bottomrule
    \end{tabular}
    \end{center}

    \textbf{问题：}
    \begin{enumerate}[(1)]
        \item 根据上表计算：准确率（Accuracy）、精确率（Precision）、召回率（Recall）、F1 分数。（保留两位小数）

        \item 在医疗诊断场景中，\textbf{漏诊}（FN）和\textbf{误诊}（FP）哪个危害更大？基于此，应该优化精确率还是召回率？

        \item F1 分数是精确率和召回率的\textbf{调和平均}而非算术平均，为什么？调和平均有什么特性使它更适合这个场景？

        \item 若将分类阈值从 0.5 调低为 0.3（更容易判定为正类），精确率和召回率会如何变化？（定性描述，不需要数值计算）
    \end{enumerate}
    {\color{gray}\small\textit{来源溯源：[文档第 4 章 4.7.2 分类模型评价指标]}}
\end{problembox}

\begin{solutionbox}[参考答案与详细解析]
    \textbf{(1) 各指标计算：}

    总样本数 $= 40+10+20+130 = 200$

    \[
    \text{Accuracy} = \frac{TP+TN}{n} = \frac{40+130}{200} = \frac{170}{200} = \mathbf{0.85}
    \]
    \[
    \text{Precision} = \frac{TP}{TP+FP} = \frac{40}{40+20} = \frac{40}{60} \approx \mathbf{0.67}
    \]
    \[
    \text{Recall} = \frac{TP}{TP+FN} = \frac{40}{40+10} = \frac{40}{50} = \mathbf{0.80}
    \]
    \[
    \text{F1} = \frac{2 \times \text{P} \times \text{R}}{\text{P}+\text{R}} = \frac{2 \times 0.67 \times 0.80}{0.67+0.80} = \frac{1.072}{1.47} \approx \mathbf{0.73}
    \]

    \textbf{(2) 医疗场景中优化目标：}

    漏诊（FN，患病却被判为未患病）危害更大——患者错过治疗时机，可能危及生命。误诊（FP，未患病被判为患病）虽然不好，但可通过进一步检查纠正，危害相对可控。因此应\textbf{优化召回率（Recall）}，尽量找出所有真实患者，宁可多些误报，不可漏掉真正患者。

    \textbf{(3) 调和平均 vs 算术平均：}

    调和平均 $H = \frac{2PR}{P+R}$ 的特性：当 P 和 R 差距很大时，调和平均会被\textbf{较小的那个拉低}（比算术平均更低）。例如 P=0.99, R=0.01，算术平均 $= 0.5$（看起来还不错），调和平均 $\approx 0.02$（真实反映了模型几乎失效）。这使 F1 不能通过"一好一坏"来"蒙混过关"，必须两者都高才能有高 F1——更准确地反映模型在精确率和召回率上的综合表现。

    \textbf{(4) 阈值降低的影响：}

    阈值降低 → 更多样本被预测为正类 → TP 增加（更多真正患者被识别）、FP 也增加（更多健康人被误判为患病）、FN 减少。结果：\textbf{召回率上升}（分母 TP+FN 中，FN 减少），\textbf{精确率下降}（分母 TP+FP 中，FP 增加）。精确率-召回率之间存在权衡（Precision-Recall Trade-off）。
\end{solutionbox}

%% ============================================================
\Needspace{5\baselineskip}
\begin{problembox}[Q4-08\quad ROC 曲线与 AUC：原理、绘制与解读]
    \textbf{题目描述：}\\
    ROC 曲线（Receiver Operating Characteristic）是评估二分类器性能的重要工具，AUC（Area Under Curve）是其综合指标。

    \vspace{0.2cm}
    \textbf{问题：}
    \begin{enumerate}[(1)]
        \item ROC 曲线的横轴和纵轴分别是什么？写出两者的公式，并说明"完美分类器"和"随机猜测分类器"在 ROC 图上各对应什么形状。

        \item AUC 的取值范围是什么？AUC $= 0.5$、AUC $= 1.0$、AUC $< 0.5$ 各代表什么含义？

        \item ROC 曲线是如何随分类阈值变化而描绘出来的？（即从阈值最高到最低，点如何从 ROC 图的左下角移动到右上角）

        \item 在类别严重不平衡（正负比例 1:99）的情况下，ROC-AUC 和 PR-AUC（Precision-Recall 曲线下面积）哪个更能反映模型真实性能？为什么？
    \end{enumerate}
    {\color{gray}\small\textit{来源溯源：[文档第 4 章 4.7.2 ROC曲线 / AUC值]}}
\end{problembox}

\begin{solutionbox}[参考答案与详细解析]
    \textbf{(1) ROC 轴与特殊分类器：}

    \begin{itemize}
        \item 横轴：\textbf{FPR（False Positive Rate）}$= \dfrac{FP}{FP+TN}$（假阳性率，负样本中被错判为正的比例）
        \item 纵轴：\textbf{TPR（True Positive Rate/Recall）}$= \dfrac{TP}{TP+FN}$（真阳性率，正样本中被正确识别的比例）
        \item \textbf{完美分类器}：ROC 曲线经过 $(0,0) \to (0,1) \to (1,1)$，即左上角折线，AUC $= 1.0$。
        \item \textbf{随机猜测}：ROC 曲线是从 $(0,0)$ 到 $(1,1)$ 的对角线，AUC $= 0.5$。
    \end{itemize}

    \textbf{(2) AUC 含义：}

    AUC $\in [0, 1]$：AUC $= 1.0$：完美分类，所有正负样本完全可分；AUC $= 0.5$：等同于随机猜测，模型无判别力；AUC $< 0.5$：比随机猜测还差（如果取反预测反而好于 0.5）。直觉解释：AUC 等于\textbf{随机取一个正样本和一个负样本，正样本预测概率高于负样本的概率}——直接量化了模型的"排序能力"。

    \textbf{(3) ROC 曲线的绘制原理：}

    从\textbf{阈值最高（仅最高置信度的正样本被预测为正）开始}，此时 TP 和 FP 都接近 0，点在 $(0,0)$ 附近。随着阈值降低，越来越多的样本被预测为正：每遇到一个真正正样本（TP++），点沿\textbf{纵轴方向上移}；每遇到一个真正负样本（FP++），点沿\textbf{横轴方向右移}。最终阈值降到最低时，所有样本都被预测为正，TP 和 FP 均最大，点到达 $(1,1)$。

    \textbf{(4) 类别不平衡下的指标选择：}

    \begin{keybox}[严重不平衡时应用 PR-AUC]
        ROC-AUC 在严重不平衡时\textbf{过于乐观}：分母 TN（大量负样本）使 FPR 很小，即使 FP 数量不少，FPR 也可能看起来很低，给出虚假高 AUC。而 PR 曲线的精确率分母是 TP+FP（预测为正的样本），不受大量负样本影响，能更真实地反映模型在少数正类上的表现——PR-AUC 更适合不平衡数据的评估。
    \end{keybox}
\end{solutionbox}

\section{本章小结}
\begin{itemize}
    \item \textbf{损失函数决定优化方向}：回归用 MSE/MAE；二分类用对数损失；多分类用交叉熵。有异常值用 MAE，否则 MSE 更常见。
    \item \textbf{欠拟合/过拟合的诊断核心}：看训练误差和验证误差的关系——都高是欠拟合，训练低验证高是过拟合。
    \item \textbf{正则化三兄弟}：L1 稀疏（特征选择）；L2 均匀压缩（稳定）；ElasticNet 兼顾两者。
    \item \textbf{梯度下降实战}：深度学习用 MBGD（mini-batch）+ Adam；学习率太大发散，太小收敛慢。
    \item \textbf{分类指标的权衡}：类别平衡且无特殊关注→准确率；不平衡或关注少数类→F1/AUC-PR；需要排序能力→ROC-AUC。
\end{itemize}
\end{document}
